\section{Discussion}\label{sec:discussion}

\subsection{Principal findings}
This study compared five foundation models with clinical, radiomic and tumour-centred models for overall survival in NSCLC, within one protocol and on the same patients and folds. Four findings emerge.

First, frozen whole-volume representations contained prognostic information, but not more than the established predictors could be shown to provide. COLIPRI had the highest internal discrimination and exceeded \model{Stage+Volume} and radiomics by 0.026 to 0.029 in point estimate with the primary head and by 0.043 to 0.046 with PCA32, but no difference was significant after correction. Among the blocks added to the clinical covariates, only tumour volume improved discrimination significantly; no encoder was shown to add to \model{Clinical+Volume}. COLIPRI did add to the model that already contained radiomics, by 0.047, so it carried information that radiomics did not. Second, more training did not help: neither trainable heads nor LoRA or partial fine-tuning improved on a linear Cox head on the frozen embedding, and the encoders were no more label-efficient than the baselines. Third, the survival predictions depended little on the tumour. Hiding the tumour inside the encoder changed discrimination by at most 0.013, although for COLIPRI it changed the embedding substantially and halved the decodable tumour volume. Fourth, risk ordering transported to an external surgical cohort for TANGERINE, DINOv3 and COLIPRI, whereas their incremental value over clinical models could not be resolved and their absolute risks were poorly calibrated.

\subsection{Comparison with previous work}
The absence of a benefit from adaptation agrees with the two studies that compared frozen and adapted encoders for prognosis: frozen lesion-encoder features outperformed fine-tuned variants,\cite{pai2024} and LoRA and full fine-tuning of CT-CLIP gave no advantage over the frozen encoder.\cite{allgower2026} Our results extend this to five encoders, to label budgets from 10\% to 100\% and to a paired statistical comparison. They also suggest one reason: in about 330 patients per development set, the trainable models selected very different numbers of epochs across folds, an indication that the signal available for optimization is weak relative to its noise. The fixed adaptation configuration did not train at all for CT-CLIP, so a configuration transferred between encoders is not a neutral choice.

The competitiveness of radiomics and of simple clinical predictors is consistent with recent benchmarks on LUNG1, in which handcrafted features matched foundation-model features\cite{neukirch2026} and frozen lesion-centred embeddings discriminated only modestly.\cite{aerts2025tib} Those benchmarks used binary two-year endpoints and frozen heads; with a censoring-aware endpoint, trainable and adapted models, and paired equivalence testing, the conclusion is the same. The best internal estimates (0.59 to 0.63) also lay within the range reported for prognostic models in this cohort (Section~\ref{sec:cohorts}).\cite{aerts2014}

Our tumour analyses reach a different conclusion from earlier work on tumour dependence. In the network of Hosny and colleagues, removing the context around the tumour reduced discrimination,\cite{hosny2018} and occlusion saliency has localized to the lesion for some encoders.\cite{aerts2025tib} Those analyses studied tumour-centred models or perturbed the input image. Intervening inside frozen whole-volume encoders, with a matched control, showed that their survival predictions were largely carried by the rest of the scan.

\subsection{What the representations encode}
The probing and masking results suggest how a whole-volume embedding can predict survival without relying on the tumour. COLIPRI encoded tumour size, and it was the only encoder whose embedding and fitted model responded to hiding the tumour; yet losing half of the decodable volume cost little discrimination, so the remaining information was largely redundant with it. The other encoders encoded tumour volume weakly but sex strongly and, for TANGERINE and DINOv3, age, consistent with representations dominated by patient-level anatomy rather than by the lesion. This may explain why the same encoders transported externally. In an early-stage surgical cohort, many deaths are likely not caused by the cancer, and representations of age and general health could rank patients as well as stage and volume, which vary little in that population; this was the expectation recorded in the external protocol. TANGERINE and DINOv3, which encoded age and sex but little tumour information, transported best. These explanations are post hoc. Decodability shows that a variable is encoded, not that the survival prediction uses it, and the analyses do not identify a biological mechanism.

\subsection{Limits of the available input}
Discrimination converged across model families. Internally, the clinical covariates, tumour volume, radiomics, the tumour-centred networks and the five encoders all lay between about 0.53 and 0.64, and neither trainable heads, encoder adaptation nor additional training patients moved the leading encoders beyond this range; externally, the three stronger encoders lay within about 0.02 of \model{Clinical+Volume}. When representations that differ this much in origin and capacity reach a similar level, the limit is more plausibly set by what the input contains than by how it is modelled. Overall survival depends on factors that pretreatment CT and the recorded clinical variables reflect only indirectly, such as performance status, comorbidity, lung function and the treatment actually delivered; neither cohort records performance status or comorbidity, and LUNG1 does not record the treatment received. The encoders appear to recover part of this information through proxies: they encoded age and sex and drew on the patient rather than on the tumour, which may explain why they matched, but did not exceed, the clinical models. Two considerations temper this interpretation. Recording these variables is not sufficient by itself: a model developed at the institution that provided LUNG1, which included performance status, lung function and delivered dose, reached a concordance of 0.62,\cite{oberije2015} close to the level observed here, so part of this limit may reflect the intrinsic unpredictability of overall survival, which includes deaths unrelated to the cancer and events that depend on care after baseline. And with a few hundred patients, differences below the minimum detectable difference cannot be excluded, so the convergence holds at the resolution of this study and does not show that larger encoders, or adaptation tuned for each encoder, could not help. Further progress therefore seems more likely to come from cohorts that record the patient and the treatment alongside the image, and from larger samples, than from larger or more extensively adapted encoders.

\subsection{Implications}
For clinical use, the results do not support replacing or augmenting clinical predictors with frozen foundation-model embeddings: no incremental value over the clinical covariates and tumour volume was demonstrated internally or externally, and absolute risks would require recalibration in every new population. For research, they argue for three practices. Foundation-model survival studies should report paired comparisons against a parsimonious clinical model and against radiomics, with equivalence margins and minimum detectable differences, because differences of a few hundredths are the scale of interest and are below the detection limit of cohorts of a few hundred patients. Adaptation choices should be selected without the test folds, and their stability checked per encoder. Finally, claims that a model ``looks at the tumour'' need interventions with matched controls, because a change in the representation need not change the prediction.

\section{Limitations}\label{sec:limitations}
The main limitation of this study is its size. The differences between models were a few hundredths of Uno C, often smaller than either cohort could reliably detect: the minimum detectable difference, the smallest difference detectable with 80\% power, was about 0.02 to 0.07 internally and, with 33 deaths by 36 months, 0.09 to 0.20 externally. Most contrasts were therefore inconclusive rather than negative, and the external analysis can neither establish nor exclude incremental value. The uncertainty we report is also likely to be underestimated. All internal analyses used a single partition into five outer folds, and the bootstrap intervals condition on the fitted models and on that partition, so they do not include the variability from refitting on different patients or from a different fold assignment.

A second group of limitations concerns how the development cohort was used. LUNG1 served for every internal analysis. The PCA32 head was added after the PCA16 comparisons, the external LoRA analysis after the external cohort had been examined, and the label-budget, masking and probing analyses are exploratory, with intervals not adjusted for multiplicity. The adaptation itself was limited in scope. A single configuration was transferred to all encoders without tuning; it failed for CT-CLIP, which was moreover adapted at its visual projection rather than at an encoder block, so its results do not show what a tuned configuration would achieve, and full fine-tuning was not examined. The clinical reference was restricted to the covariates shared by both cohorts, and neither cohort records performance status or comorbidity, so the incremental value of the encoders was assessed against a weaker clinical model than is available in practice, and the interpretation that the available input limits discrimination could not be tested directly.

The comparisons between pipelines have their own limits. Each encoder kept its published preprocessing and global pooling, so the differences combine pretraining, preprocessing, field of view and pooling, and cannot be attributed to the pretraining objective; the tumour-centred comparators, in turn, received the tumour location that the whole-volume encoders did not. The comparators were not all exact reproductions: the Hosny network, reimplemented from the published weights, did not fully reproduce the published patient-level predictions on LUNG1, and the radiomic model is a practical baseline rather than the original LUNG1 signature. The tumour interventions also require caution. Hidden tokens contain the tissue around the tumour within their footprint, the masked inputs lie outside the distribution on which the encoders were trained, and regional pooling locates information by token position rather than by tissue of origin, because tokens exchange information inside the encoder. Robustness to masking therefore does not show that the tumour carries no prognostic information.

Finally, the external cohort differs from LUNG1 in stage, treatment, scanners and the source of TNM, which was pathological rather than clinical. Survival was measured from the CT date because treatment dates were unavailable, and the segmentation labels required manual verification. In LUNG1, the legacy codes T5 and N4 were merged by an analytic convention whose clinical meaning is uncertain. Some encoders were pretrained on public chest CT collections, and the overlap between their pretraining data and either cohort has not been fully audited.

\section{Conclusions}\label{sec:conclusions}
Frozen CT foundation models captured prognostic information in NSCLC comparable to, but not demonstrably greater than, stage, tumour volume and radiomics; trainable heads and encoder adaptation did not improve on a linear head; and the predictions depended little on the tumour itself. Their risk ordering transported to an external cohort, but their incremental clinical value remains unresolved and their absolute risks require recalibration. Because discrimination converged across clinical, radiomic and foundation models, the input available in public cohorts, rather than the choice of model, appears to limit prognostic performance. Larger cohorts that record the patient and the treatment, prespecified adaptation choices and evaluation of incremental value over clinical models are needed before these representations can be considered for prognostic use.
